% FILE: 3-poster/parts/colors-values.tex
% ==================================================================================================
% LAYER 1 of the colour system: WHAT EACH TOKEN IS. The only place a colour may be DEFINED.
% Loaded first by parts/colors.tex, which then applies the scheme's overrides and finally binds the
% roles. See that file for the layering and for why the order is not interchangeable.
%
% >>> NO TOKEN HERE MAY BE NAMED FOR ITS HUE <<< A name states a USE (`inputbright', `comp1') or a
% rung of a ladder, never a colour word. A hue-named token silently becomes a lie the moment a scheme
% retunes it -- `accentturq' holding a grey is worse than no documentation at all.
%
% >>> THIS FILE IS NOT WHERE A COLOUR IS ASSIGNED TO AN ELEMENT <<< A call site never names a token
% from this file for a CHROME or NEUTRAL use; it names a role (parts/colors-roles.tex). Figures still
% name semantic tokens directly, because a semantic token IS its own role: one concept, one colour.
%
% THREE INDEPENDENT PALETTES. Each token belongs to exactly one, so any palette can be retuned alone:
%   1. Neutral foundation -- achromatic substrate: background, surfaces, borders, text, plot axes.
%                            Shared by chrome and figures; never an accent, never a concept.
%   2. Design chrome      -- the poster's own identity colour (+ a reserved complement): headings,
%                            boxes, navigation, QR tags, badges, and both kinds of pointer AWAY from
%                            the poster (\litcite, \seebottom -- one colour, one kind of mark).
%   3. Semantic grammar   -- one concept, one colour. Figures only.
% ONE LICENSED EXCEPTION: `modelnet' encodes a concept (the poster's own model) in a colour close to
% the chrome, in the model comparison only (see 3c). Nowhere else may chrome mean something. It is a
% SEPARATE token, not an alias, so a scheme that neutralises the chrome leaves the model comparison
% standing -- which is the point of that scheme.
%
% ADDING OR RETUNING A TOKEN -- the acceptance test, in this order:
%   (a) dE >= 0.15 in OKLab against every token a reader must tell it apart from. Below ~0.15 two
%       colours are not distinct CATEGORIES, whatever they look like side by side.
%   (b) >= 3:1 contrast on `posterbg' if it is a graphical mark. Fills and washes are exempt and must
%       carry a saturated stroke instead.
%   (c) the ladder limit (3a) if it belongs to a pole.
% Values are sRGB approximations of an OKLCH specification: the lightness ORDER within a hue is what
% matters, not exact colorimetry.

% ----- 1. Neutral foundation ----------------------------------------------------------------------
\definecolor{posterbg}{HTML}{F1F3F3}      % L=0.970  background
\definecolor{surface}{HTML}{FBFCFC}       % L=0.995  elevated panels (shadow separates from bg)
% THE TWO QUESTION CARDS' GROUND, read ONLY in the OUTLINED treatment (\questioninvertedfalse); the
% inverted treatment grounds them in the accent instead. It exists so a scheme can tint the page's two
% HINGES without tinting every card, which moving `surface' would. A VALUE and not a role, because the
% role layer loads after the scheme and a scheme may not bind a role.
% >>> KEEP IT EQUAL TO `surface' UNLESS THE TINT IS THE POINT <<< A scheme that retunes `surface' and
% leaves this behind grounds the question cards in the OLD surface, and nothing reports it.
\definecolor{questionwash}{HTML}{FBFCFC}  % L=0.995  == surface, until a scheme tints it
\definecolor{bordersubtle}{HTML}{DBDFDF}  % L=0.900  decorative outlines, dividers, table lines
\definecolor{borderfunc}{HTML}{8C9698}    % L<=0.640 focus rings, emphasis boundaries, plot axes
\definecolor{texttertiary}{HTML}{6E797B}  % L=0.530  captions, minor metadata, faint hints
\definecolor{textsecondary}{HTML}{586366} % L=0.450  subtitles, supporting text
\definecolor{textprimary}{HTML}{1B2224}   % L=0.150  titles, labels, body, key values

% ----- 2. Design chrome ---------------------------------------------------------------------------
% THE POLE A SCHEME RETUNES. Everything chrome on the page reaches these three through a role, so
% overriding `chromecore' alone repaints the whole identity layer and nothing else.
\definecolor{chromecore}{HTML}{12A89E}  % core -- the identity colour (H 187)
\definecolor{chromedeep}{HTML}{0C7A72}  % deep -- sub-headers, contact line, badges (carries white text)
\definecolor{chromealt}{HTML}{F26B5B}   % complement -- RESERVED, chrome only (H 12)
% The medallion network's body wash (figures/network-medallion.tex). CHROME, NOT A CONCEPT: it says
% "one network", which is the container the neurons sit in, not a claim about it -- so it belongs to
% this palette and not to the semantic grammar, and it must never be given a meaning. A NAMED token
% rather than an `chromecore!12' at the call site, because a hand-mixed tint is exactly what this file
% exists to prevent. NEARLY INVISIBLE BY DESIGN, at L=0.982 -- a hair off the card's `surface' (0.995)
% and barely tinted: the disk's job is to give the neurons a ground, and a wash strong enough to be
% noticed competes with the palest neurons, which are themselves within a few percent of white. What
% marks the network is its BOUNDARY (4.5 pt of accent), not its fill.
\definecolor{netfill}{HTML}{F7FAFA}       % network body wash -- FILL: carries no meaning

% ==================================================================================================
% 3. SEMANTIC GRAMMAR -- one concept, one colour. Figures only.
% ==================================================================================================
%
% ----- 3a. The two transformations: a blue pole (H 258) and a magenta pole (H 330) -----------------
% THE COMPONENTS OWN RED/GREEN; THE TRANSFORMATIONS OWN BLUE/MAGENTA. Keep the two in DISJOINT hue
% regions: the components name WHAT IS COMBINED, the transformations name HOW THE NETWORK ACTS, and two
% concepts of different kinds must not sit in adjacent hues. Measured constraint -- an amplification
% hue near component 2 lands at dE 0.107 from it, i.e. indistinguishable as categories in the figures
% where both appear. Do not move either pole toward the other.
% Blue reads exogenous and directional, magenta endogenous and self-reinforcing, which is the model:
% the input arrives from outside, the recurrence feeds back on itself.
%
% ONE LADDER PER POLE: each transformation is a constant-hue tint of its own SOURCE, so the projection
% is visibly the input's operation and the amplification the recurrence's.
%   deep    card chrome (carries white text)   core    the source concept (input I; connectivity J)
%   bright  the TRANSFORMATION (xi_k ; A_k)    light   a second pattern within the pole
%   pale    fills and washes
%
% >>> THE LADDER LIMIT -- IT BINDS <<<
%   (1) Adjacent rungs sit at dE ~0.09: NOT distinct. Within one figure use at most THREE levels of a
%       hue and SKIP a rung between them (deep, bright, pale). Never two adjacent rungs together.
%   (2) The two pale fills are dE 0.041 apart, so PALE CANNOT ENCODE THE POLE: always pair a pale fill
%       with a saturated stroke of its own family.
%   (3) A fourth distinction goes to a NON-COLOUR channel. A pattern's index (first or second) is an
%       index, not a concept: it belongs in a subscript or a line style, never a fourth tint.
% figures/approach-steps.tex is the figure this limit was written for, and it still strains it.
%
% --- blue pole: the external input, and what it does ---
\definecolor{inputdeep}{HTML}{0F4EA3}   % deep   -- input pattern 1 (I^(1)); projection card chrome
\definecolor{inputcore}{HTML}{2368C9}     % core   -- the external input itself (I, u)
\definecolor{inputbright}{HTML}{3483F3}      % bright -- THE FEEDFORWARD PROJECTION (xi_k, Sigma^(nI))
\definecolor{inputlight}{HTML}{74A6EF}   % light  -- input pattern 2 (I^(2))
\definecolor{inputpale}{HTML}{D7E6FC}    % pale   -- input pooling Sigma^(II)   [FILL: needs a stroke]
% --- magenta pole: the recurrence, and what it does ---
\definecolor{recurdeep}{HTML}{7D2A79}    % deep   -- recurrent pattern 1 (m_1); amplification chrome
\definecolor{recurcore}{HTML}{9E4098}      % core   -- the recurrent connectivity itself (J, Sigma^(nm))
\definecolor{recurbright}{HTML}{C154BA}        % bright -- THE RECURRENT AMPLIFICATION (A_k)
\definecolor{recurlight}{HTML}{CF88C8}    % light  -- recurrent pattern 2 (m_2)
\definecolor{recurpale}{HTML}{F3DDF0}    % pale   -- recurrent pooling Sigma^(mm) [FILL: needs a stroke]

% ----- 3b. The components, the response, the geometry ----------------------------------------------
\definecolor{comp1}{HTML}{D6453C}       % component response 1 (kappa_1)   H 28
\definecolor{comp2}{HTML}{4F9A3F}     % component response 2 (kappa_2)   H 140
\definecolor{mixresponse}{HTML}{F2A30C}      % THE RESPONSE (H 74) -- mixture kappa_mix, network state x,
                                          % reduced state kappa: ONE concept, ONE token. Do not split
                                          % it; a second token for the same object would land between
                                          % the two poles and read as a blend of them.
\definecolor{compplane}{HTML}{B7BFC1}      % components plane / proportional reference
\definecolor{outplaneannot}{HTML}{566063}  % out-of-plane annotations (bias arc, in-plane foot marker)
\definecolor{outplane}{HTML}{8A5524} % orthogonal component (eta_perp) -- solid brown, all figures

% ----- 3c. The two models: chromatic vs achromatic, not a third hue pair ---------------------------
% THE HUE CIRCLE IS FULL -- do not spend a new semantic hue here. Occupied: H 28 (component 1),
% 60 (eta_perp), 74 (the response), 140 (component 2), 187 (chrome), 258 (feedforward),
% 330 (amplification). The widest remaining gap spans 72 deg (+/-36 deg of clearance) against the
% ~45 deg a categorical hue needs at this chroma, so anything forced in lands ~dE 0.13 from BOTH
% transformations. An achromatic PAIR fails too: its light member sits dE 0.034 from borderfunc at
% 2.4:1 on the background, under the 3:1 mark floor. Hence chromatic/achromatic, doubled in LINE STYLE
% so the distinction survives a grayscale print.
\definecolor{modelnet}{HTML}{00948B}     % THE NETWORK MODEL (this work): SOLID, filled markers. The
                                          % licensed chrome exception; model-comparison band only.
\definecolor{modeldn}{HTML}{474E51}       % DIVISIVE NORMALIZATION (rival): DASHED, hollow markers.
                                          % A mechanism earns a hue; a descriptive fit does not.
% modeldn sits dE 0.063 from outplaneannot. THEY NEVER APPEAR IN ONE FIGURE -- re-check if either moves.

% ----- END OF THE VALUE LAYER ---------------------------------------------------------------------
% Every \colorlet that USED to close this file now lives in parts/colors-roles.tex, and that move is
% load-bearing rather than cosmetic: \colorlet COPIES A VALUE AT THE MOMENT IT IS READ. An alias left
% here would snapshot the base value and then survive a scheme override untouched, repainting half the
% page and leaving the other half in the base's accent, with nothing in the log to say so.
